\chapter{Quadratura de Gauss-Lobatto}

\indent Uma das formas mais clássicas de se aproximar uma integral de funções de uma variável é a regra do trapézio. O resultado de uma integral definida de uma função $f(x)$ entre dois pontos $x_1$ e $x_2$ é equivalente à área sob o gráfico de $f(x)$ entre esses pontos. A regra do trapézio consiste em calcular a área da integral por meio da soma da área de diversos trapézios com vértices equivalentes aos pontos $x_{i}$ e $x_{i+1}$ e à avaliação da função $f(x)$ nesses pontos.\\
\indent Segundo \citet{burden2015numerical}, a técnica de quadratura busca otimizar a escolha dos pontos de avaliação da integral em vez de utilizar distâncias igualmente espaçadas. Para isso, são determinados pontos $x_1, x_2, \dots, x_N$ no intervalo $[a,b]$ e coeficientes $c_1, c_2, \dots, c_N$ são escolhidos para minimizar o erro esperado, de modo que: 
\begin{equation}
  \int_a^b f(x) \, dx \approx \sum_{i=1}^N c_i f(x_i).
\end{equation}
\noindent Uma forma de determinar o valor dos pontos e coeficientes é utilizando os polinômios de Legendre.\\
\indent Os polinômios de Legendre ortogonais em $[-1, 1]$ apresentam propriedades importantes. Na quadratura de Gauss-Legendre padrão, os pontos correspondem às raízes do polinômio de Legendre de grau $N$, e os coeficientes são obtidos de:
\begin{equation}
  w_i = \int_{-1}^1 L_i(x) \, dx,
\end{equation}
\noindent onde $L_i(x)$ é o $i$-ésimo polinômio base de Lagrange.\\
\indent \citet{kopriva2009implementing} comenta que, na quadratura de Gauss-Lobatto, os pontos extremos devem ser obrigatoriamente $-1$ e $1$, e os pontos internos correspondem às raízes da derivada de um polinômio de Legendre de grau $N-1$. Os coeficientes dessa quadratura são determinados analogamente por meio de:
\begin{equation}
  w_j = \int_{-1}^1 L_j(x) \, dx.
\end{equation}
\indent A tabela abaixo apresenta os pontos e os pesos da quadratura de Gauss-Lobatto para ordens de aproximação crescentes:

\begin{table}[htpb]
\centering
\begin{tabular}{|c|c|r|r|}
\hline
\rowcolor[HTML]{C0C0C0} 
{\color[HTML]{333333} \textbf{N}} & 
{\color[HTML]{333333} \textbf{Índice ($i$)}} & 
{\color[HTML]{333333} \textbf{Ponto ($x_i$)}} & 
{\color[HTML]{333333} \textbf{Peso ($w_i$)}} \\ \hline
\multirow{2}{*}{2} 
& 0 & -1,000000 & 1,000000 \\ \cline{2-4}
& 1 & 1,000000 & 1,000000 \\ \hline
\multirow{3}{*}{3} 
& 0 & -1,000000 & 0,333333 \\ \cline{2-4}
& 1 & 0,000000 & 1,333333 \\ \cline{2-4}
& 2 & 1,000000 & 0,333333 \\ \hline
\multirow{4}{*}{4} 
& 0 & -1,000000 & 0,166667 \\ \cline{2-4}
& 1 & -0,577350 & 0,833333 \\ \cline{2-4}
& 2 & 0,577350 & 0,833333 \\ \cline{2-4}
& 3 & 1,000000 & 0,166667 \\ \hline
\multirow{5}{*}{5} 
& 0 & -1,000000 & 0,100000 \\ \cline{2-4}
& 1 & -0,654654 & 0,544444 \\ \cline{2-4}
& 2 &  0,000000 & 0,711111 \\ \cline{2-4}
& 3 &  0,654654 & 0,544444 \\ \cline{2-4}
& 4 &  1,000000 & 0,100000 \\ \hline
\multirow{6}{*}{6} 
& 0 & -1,000000 & 0,066667 \\ \cline{2-4}
& 1 & -0,765055 & 0,378475 \\ \cline{2-4}
& 2 & -0,285232 & 0,554858 \\ \cline{2-4}
& 3 &  0,285232 & 0,554858 \\ \cline{2-4}
& 4 &  0,765055 & 0,378475 \\ \cline{2-4}
& 5 &  1,000000 & 0,066667 \\ \hline
\multirow{7}{*}{7} 
& 0 & -1,000000 & 0,047619 \\ \cline{2-4}
& 1 & -0,830224 & 0,276826 \\ \cline{2-4}
& 2 & -0,468849 & 0,431745 \\ \cline{2-4}
& 3 &  0,000000 & 0,487619 \\ \cline{2-4}
& 4 &  0,468849 & 0,431745 \\ \cline{2-4}
& 5 &  0,830224 & 0,276826 \\ \cline{2-4}
& 6 &  1,000000 & 0,047619 \\ \hline
\multirow{8}{*}{8} 
& 0 & -1,000000 & 0,035714 \\ \cline{2-4}
& 1 & -0,871740 & 0,210704 \\ \cline{2-4}
& 2 & -0,591700 & 0,341123 \\ \cline{2-4}
& 3 & -0,209299 & 0,412459 \\ \cline{2-4}
& 4 &  0,209299 & 0,412459 \\ \cline{2-4}
& 5 &  0,591700 & 0,341123 \\ \cline{2-4}
& 6 &  0,871740 & 0,210704 \\ \cline{2-4}
& 7 &  1,000000 & 0,035714 \\ \hline
\end{tabular}
\end{table}